\chapter{Bogoliubov--de Gennes Formalism}

The Jordan--Wigner transformation maps several spin-$1/2$ chains to quadratic
fermion Hamiltonians.  If the fermion number is conserved, Fourier
transformation already diagonalizes the problem into one-mode momentum blocks. If
pairing terms are present, however, the Hamiltonian conserves only fermion
parity, and the elementary block is no longer a single mode $k$ but a paired
sector $(k,-k)$.  The Bogoliubov--de Gennes (BdG) formalism is the systematic
language for this situation.

The purpose of this chapter is to fix the algebraic conventions used later in
the XY chain, the transverse-field Ising chain, and the Kitaev chain.  We do not
yet perform the full Bogoliubov rotation; that will be done in the next chapter.
Here we explain how a quadratic pairing Hamiltonian is encoded as a matrix in
Nambu space, why the resulting matrix has a built-in particle--hole redundancy,
and how this Nambu-space structure differs from an ordinary two-band Bloch
Hamiltonian such as the SSH model.

\section{Real-space quadratic fermion Hamiltonians}

Let $c_j,c_j^\dagger$ be canonical spinless fermions on a chain of length $L$,
obeying
\begin{equation}
    \{c_i,c_j^\dagger\}=\delta_{ij},
    \qquad
    \{c_i,c_j\}=\{c_i^\dagger,c_j^\dagger\}=0.
\end{equation}
The most general quadratic Hamiltonian preserving fermion parity is
\begin{equation}
\label{eq:bdg-real-space-quadratic-hamiltonian}
    H
    =
    \sum_{i,j=1}^{L} h_{ij}c_i^\dagger c_j
    +\frac12\sum_{i,j=1}^{L}
    \left(
        \Delta_{ij}c_i^\dagger c_j^\dagger
        +\Delta_{ij}^*c_j c_i
    \right)
    +E_{\rm cl}.
\end{equation}
Hermiticity requires
\begin{equation}
    h^\dagger=h.
\end{equation}
Fermi statistics require the pairing matrix to be antisymmetric,
\begin{equation}
\label{eq:bdg-pairing-antisymmetry-real-space}
    \Delta^T=-\Delta,
\end{equation}
because only the antisymmetric part of $\Delta_{ij}c_i^\dagger c_j^\dagger$
survives.

It is useful to introduce the real-space Nambu spinor
\begin{equation}
\label{eq:bdg-real-space-nambu-spinor}
    \Upsilon
    =
    \begin{pmatrix}
        c_1\\
        \vdots\\
        c_L\\
        c_1^\dagger\\
        \vdots\\
        c_L^\dagger
    \end{pmatrix},
    \qquad
    \Upsilon^\dagger
    =
    \begin{pmatrix}
        c_1^\dagger&\cdots&c_L^\dagger&c_1&\cdots&c_L
    \end{pmatrix}.
\end{equation}
Then \cref{eq:bdg-real-space-quadratic-hamiltonian} can be written as
\begin{equation}
\label{eq:bdg-real-space-matrix-form}
    H
    =
    \frac12\Upsilon^\dagger\mathbb{H}_{\rm BdG}\Upsilon
    +E_{\rm const},
\end{equation}
where
\begin{equation}
\label{eq:bdg-real-space-matrix}
    \mathbb{H}_{\rm BdG}
    =
    \begin{pmatrix}
        h & \Delta\\
        -\Delta^* & -h^T
    \end{pmatrix}.
\end{equation}
The constant differs from the normal-ordering constant in
\cref{eq:bdg-real-space-quadratic-hamiltonian}; explicitly,
\begin{equation}
    E_{\rm const}=E_{\rm cl}+\frac12\Tr h.
\end{equation}
This shift has no effect on eigenvectors or quasiparticle energies, but it must
be retained when computing the absolute ground-state energy.

\begin{remark}[Why the factor of $1/2$ appears]
The Nambu spinor contains both particles and holes.  Therefore the degrees of
freedom are doubled.  The prefactor $1/2$ in
\cref{eq:bdg-real-space-matrix-form} compensates this doubling.  Without this
factor the normal hopping term and the pairing term would be counted twice.
\end{remark}

\section{Intrinsic particle--hole redundancy}

The BdG matrix is not an arbitrary Hermitian $2L\times 2L$ matrix.  It satisfies
an intrinsic particle--hole constraint.  Let
\begin{equation}
    \tau_x=
    \begin{pmatrix}
        0&\id_L\\
        \id_L&0
    \end{pmatrix},
    \qquad
    \mathcal{C}=\tau_x K,
\end{equation}
where $K$ denotes complex conjugation.  Then
\begin{equation}
\label{eq:bdg-real-space-phs}
    \tau_x\mathbb{H}_{\rm BdG}^*\tau_x
    =
    -\mathbb{H}_{\rm BdG},
    \qquad
    \text{equivalently}\qquad
    \mathcal{C}\mathbb{H}_{\rm BdG}\mathcal{C}^{-1}
    =
    -\mathbb{H}_{\rm BdG}.
\end{equation}
Thus if
\begin{equation}
    \mathbb{H}_{\rm BdG}\Phi=E\Phi,
\end{equation}
then
\begin{equation}
    \mathbb{H}_{\rm BdG}(\mathcal{C}\Phi)=-E(\mathcal{C}\Phi).
\end{equation}
The spectrum is symmetric about zero.  This symmetry is not an optional physical
symmetry imposed on the model; it is a redundancy of the Nambu description.

\begin{principle}[BdG redundancy principle]
A BdG Hamiltonian doubles the Hilbert-space variables by treating $c$ and
$c^\dagger$ as independent components of a Nambu spinor.  The particle--hole
constraint removes this artificial doubling: the negative-energy BdG states are
particle--hole partners of the positive-energy states, not independent physical
single-particle bands.
\end{principle}

\section{Translation-invariant BdG blocks}

For a translation-invariant chain, the Fourier convention of the previous
chapter is
\begin{equation}
\label{eq:bdg-chapter-fourier-convention}
    c_k
    =
    \frac{\ee^{-\ii\phi}}{\sqrt{L}}
    \sum_{j=1}^{L}\ee^{-\ii kj}c_j,
    \qquad
    c_j
    =
    \frac{\ee^{\ii\phi}}{\sqrt{L}}
    \sum_{k\in\mathcal{K}_\vartheta}\ee^{\ii kj}c_k,
\end{equation}
with
\begin{equation}
    \mathcal{K}_\vartheta
    =
    \left\{
        k_m=\frac{2\pi m+\vartheta}{L}
        \ \middle|\
        m=0,1,\ldots,L-1
    \right\}.
\end{equation}
The boundary twist $\vartheta$ fixes the allowed momenta, while the constant
phase $\phi$ is a Fourier-gauge convention.

Consider the momentum-space Hamiltonian
\begin{equation}
\label{eq:bdg-chapter-general-momentum-hamiltonian}
    H
    =
    \sum_{k\in\mathcal{K}_\vartheta}\xi(k)c_k^\dagger c_k
    +\frac12\sum_{k\in\mathcal{K}_\vartheta}
    \left[
        \Delta_\phi(k)c_k^\dagger c_{-k}^\dagger
        +\Delta_\phi(k)^*c_{-k}c_k
    \right]
    +E_{\rm cl}.
\end{equation}
For spinless fermions,
\begin{equation}
\label{eq:bdg-chapter-odd-gap}
    \Delta_\phi(-k)=-\Delta_\phi(k).
\end{equation}
The natural Nambu spinor for the pair $(k,-k)$ is
\begin{equation}
\label{eq:bdg-chapter-nambu-spinor}
    \Psi_k
    =
    \begin{pmatrix}
        c_k\\
        c_{-k}^\dagger
    \end{pmatrix},
    \qquad
    \Psi_k^\dagger
    =
    \begin{pmatrix}
        c_k^\dagger&c_{-k}
    \end{pmatrix}.
\end{equation}
Then
\begin{equation}
\label{eq:bdg-chapter-block-form}
    H
    =
    \frac12\sum_{k\in\mathcal{K}_\vartheta}
    \Psi_k^\dagger\mathcal{H}_{\rm BdG}(k)\Psi_k
    +E_0,
\end{equation}
where
\begin{equation}
\label{eq:bdg-chapter-block-matrix}
    \mathcal{H}_{\rm BdG}(k)
    =
    \begin{pmatrix}
        \xi(k) & \Delta_\phi(k)\\
        \Delta_\phi(k)^* & -\xi(-k)
    \end{pmatrix}.
\end{equation}
The corresponding particle--hole constraint is
\begin{equation}
\label{eq:bdg-chapter-phs-momentum}
    \tau_x\mathcal{H}_{\rm BdG}(-k)^*\tau_x
    =
    -\mathcal{H}_{\rm BdG}(k).
\end{equation}
For the nearest-neighbour chains studied later, one usually has
\begin{equation}
    \xi(-k)=\xi(k),
    \qquad
    \Delta_\phi(-k)=-\Delta_\phi(k),
\end{equation}
and the matrix simplifies to
\begin{equation}
\label{eq:bdg-chapter-even-dispersion-block}
    \mathcal{H}_{\rm BdG}(k)
    =
    \begin{pmatrix}
        \xi(k) & \Delta_\phi(k)\\
        \Delta_\phi(k)^* & -\xi(k)
    \end{pmatrix}.
\end{equation}

\section{Two-level representation}

Let $\tau_x,\tau_y,\tau_z$ denote Pauli matrices acting in Nambu space, not in the
original spin space.  For the common even-dispersion case
\cref{eq:bdg-chapter-even-dispersion-block}, the BdG block can be written as
\begin{equation}
\label{eq:bdg-chapter-d-vector-form}
    \mathcal{H}_{\rm BdG}(k)
    =
    d_x(k)\tau_x+d_y(k)\tau_y+d_z(k)\tau_z,
\end{equation}
where
\begin{equation}
\label{eq:bdg-chapter-d-vector-components}
    d_x(k)=\operatorname{Re}\Delta_\phi(k),
    \qquad
    d_y(k)=-\operatorname{Im}\Delta_\phi(k),
    \qquad
    d_z(k)=\xi(k).
\end{equation}
The BdG spectrum is therefore
\begin{equation}
\label{eq:bdg-chapter-spectrum}
    E_\pm(k)
    =
    \pm E(k),
    \qquad
    E(k)=\sqrt{\xi(k)^2+\abs{\Delta_\phi(k)}^2}
    =\sqrt{d_x(k)^2+d_y(k)^2+d_z(k)^2}.
\end{equation}
A bulk gap closes only if
\begin{equation}
\label{eq:bdg-chapter-gap-closing-condition}
    d_x(k)=d_y(k)=d_z(k)=0
\end{equation}
for some allowed momentum $k$ in the thermodynamic limit.

\begin{remark}[Meaning of the negative branch]
The negative branch $-E(k)$ is not a second independent physical species.  It is
the particle--hole partner of the positive branch $E(k)$.  The physical
quasiparticle spectrum after Bogoliubov diagonalization consists of the positive
energies $E(k)$, together with a ground-state energy shift.
\end{remark}

\section[Independent (k,-k) sectors]{Independent $(k,-k)$ sectors}

For $k\not\equiv -k$ modulo the reciprocal lattice, the four-dimensional Fock
space generated by $c_k^\dagger$ and $c_{-k}^\dagger$ splits into even and odd
fermion-parity sectors.  The even sector is
\begin{equation}
\label{eq:bdg-chapter-even-sector-basis}
    \mathcal{H}_{k,-k}^{\rm even}
    =
    \operatorname{span}\left\{
        \ket{0}_{k,-k},
        c_k^\dagger c_{-k}^\dagger\ket{0}_{k,-k}
    \right\},
\end{equation}
whereas the odd sector is
\begin{equation}
    \mathcal{H}_{k,-k}^{\rm odd}
    =
    \operatorname{span}\left\{
        c_k^\dagger\ket{0}_{k,-k},
        c_{-k}^\dagger\ket{0}_{k,-k}
    \right\}.
\end{equation}
Pairing terms connect the two states in the even sector but do not connect states
of opposite fermion parity.  Thus the BdG problem is equivalent to a two-level
problem in each even-parity $(k,-k)$ sector.

If the allowed momentum set contains self-conjugate momenta satisfying
\begin{equation}
    k\equiv -k \pmod{2\pi},
\end{equation}
then $k=0$ or $k=\pi$.  For spinless odd-parity pairing,
\begin{equation}
    \Delta_\phi(0)=0,
    \qquad
    \Delta_\phi(\pi)=0,
\end{equation}
so these modes are not paired by the superconducting term.  They contribute
ordinary number-operator pieces of the form
\begin{equation}
    \xi(k)c_k^\dagger c_k.
\end{equation}
Such self-conjugate modes are often irrelevant in the thermodynamic bulk
spectrum, but they matter in finite-size calculations and in the parity-sector
bookkeeping of Jordan--Wigner-transformed spin chains.

\section{Gauge phase in Nambu space}

The previous chapter allowed a constant phase $\phi$ in the Fourier transform.
This phase acts on momentum-space fermions as
\begin{equation}
    c_k^{(\phi)}=\ee^{-\ii\phi}c_k^{(0)}.
\end{equation}
Consequently,
\begin{equation}
\label{eq:bdg-chapter-nambu-gauge-transformation}
    \Psi_k^{(\phi)}
    =
    \ee^{-\ii\phi\tau_z}\Psi_k^{(0)}.
\end{equation}
The BdG matrix transforms as
\begin{equation}
\label{eq:bdg-chapter-bdg-gauge-transformation}
    \mathcal{H}_{\rm BdG}^{(\phi)}(k)
    =
    \ee^{-\ii\phi\tau_z}
    \mathcal{H}_{\rm BdG}^{(0)}(k)
    \ee^{\ii\phi\tau_z}.
\end{equation}
Therefore $\phi$ rotates the $\tau_x$--$\tau_y$ components of the BdG vector while
leaving $d_z(k)=\xi(k)$ and the spectrum invariant.  Equivalently, the pairing
gap transforms as
\begin{equation}
\label{eq:bdg-chapter-gap-gauge-phase}
    \Delta_\phi(k)=\ee^{-2\ii\phi}\Delta_0(k).
\end{equation}
The factor of two reflects the fact that the pairing operator carries charge
$2$ under the fermion $U(1)$ phase rotation.

This distinction is important:
\begin{equation}
    \vartheta \text{ changes the allowed momenta, whereas } \phi
    \text{ only changes the Nambu gauge.}
\end{equation}
A boundary twist $\vartheta$ can represent a physical flux through a ring.  A
constant $\phi$ is a basis convention unless it is tied to a relative phase
between distinct superconducting regions or to a spatially varying order
parameter.

\section{Example: Kitaev-chain BdG matrix}

Consider the spinless Kitaev chain
\begin{equation}
\label{eq:bdg-chapter-kitaev-real-space}
    H_{\rm K}
    =
    -\mu\sum_j\left(c_j^\dagger c_j-\frac12\right)
    -t\sum_j\left(c_j^\dagger c_{j+1}+c_{j+1}^\dagger c_j\right)
    +\sum_j
    \left(
        \Delta c_j^\dagger c_{j+1}^\dagger
        +\Delta^* c_{j+1}c_j
    \right),
\end{equation}
with
\begin{equation}
    \Delta=\abs{\Delta}\ee^{\ii\theta_\Delta}.
\end{equation}
Using the Fourier convention \cref{eq:bdg-chapter-fourier-convention}, one finds
\begin{equation}
\label{eq:bdg-chapter-kitaev-xi-gap}
    \xi(k)=-\mu-2t\cos k,
    \qquad
    \Delta_\phi(k)
    =
    2\ii\abs{\Delta}\ee^{\ii(\theta_\Delta-2\phi)}\sin k.
\end{equation}
The BdG block is
\begin{equation}
\label{eq:bdg-chapter-kitaev-bdg-block}
    \mathcal{H}_{\rm K}(k)
    =
    \begin{pmatrix}
        -\mu-2t\cos k
        &
        2\ii\abs{\Delta}\ee^{\ii(\theta_\Delta-2\phi)}\sin k
        \\
        -2\ii\abs{\Delta}\ee^{-\ii(\theta_\Delta-2\phi)}\sin k
        &
        \mu+2t\cos k
    \end{pmatrix}.
\end{equation}
The spectrum is
\begin{equation}
\label{eq:bdg-chapter-kitaev-spectrum}
    E_\pm(k)
    =
    \pm\sqrt{(\mu+2t\cos k)^2+4\abs{\Delta}^2\sin^2 k}.
\end{equation}
It is independent of the uniform phase $\theta_\Delta$ and of the Fourier-gauge
phase $\phi$.

Choosing
\begin{equation}
    \phi=\frac{\theta_\Delta}{2}
\end{equation}
removes the explicit superconducting phase from the off-diagonal matrix element:
\begin{equation}
    \Delta_\phi(k)=2\ii\abs{\Delta}\sin k.
\end{equation}
For a single uniform Kitaev chain this is only a gauge choice.  In contrast,
relative superconducting phases between different segments, Josephson junctions,
or spatial phase gradients are physical because they cannot be removed by one
constant global choice of $\phi$.

For $\phi=\theta_\Delta/2$, the BdG vector is
\begin{equation}
    \bm{d}(k)
    =
    \bigl(0,-2\abs{\Delta}\sin k,-\mu-2t\cos k\bigr),
\end{equation}
using the convention $d_y=-\operatorname{Im}\Delta_\phi$.  The gap closes at
\begin{equation}
    \sin k=0,
    \qquad
    \mu+2t\cos k=0,
\end{equation}
so the thermodynamic phase boundaries are
\begin{equation}
\label{eq:bdg-chapter-kitaev-phase-boundaries}
    \mu=-2t
    \quad (k=0),
    \qquad
    \mu=2t
    \quad (k=\pi).
\end{equation}
Equivalently, for real nonzero $t$,
\begin{equation}
    \abs{\mu}=2\abs{t}
\end{equation}
marks the bulk gap closing between the trivial and topological phases.

\section{Example: BdG form of the transverse-field XY chain}

The transverse-field XY chain later studied as a mother model has the Pauli-form
Hamiltonian
\begin{equation}
\label{eq:bdg-chapter-xy-spin-hamiltonian}
    H_{\rm XY}
    =
    -\sum_{j=1}^{L}
    \left[
        \frac{1+\gamma}{2}\sigma_j^x\sigma_{j+1}^x
        +
        \frac{1-\gamma}{2}\sigma_j^y\sigma_{j+1}^y
    \right]
    -h\sum_{j=1}^{L}\sigma_j^z,
\end{equation}
up to the boundary-sector subtleties already fixed in the Jordan--Wigner chapter.
With the convention
\begin{equation}
    n_j=c_j^\dagger c_j=\frac{1-\sigma_j^z}{2},
\end{equation}
the bulk Jordan--Wigner image is a quadratic pairing Hamiltonian.  In the common
normalization used in these notes, its BdG vector may be chosen as
\begin{equation}
\label{eq:bdg-chapter-xy-d-vector}
    \bm{d}_{\rm XY}(k)
    =
    \bigl(0,-2\gamma\sin k,2(h-\cos k)\bigr),
\end{equation}
so that
\begin{equation}
\label{eq:bdg-chapter-xy-spectrum}
    E_{\rm XY}(k)
    =
    2\sqrt{(h-\cos k)^2+\gamma^2\sin^2 k}.
\end{equation}
This form makes the common free-fermion structure transparent:
\begin{equation}
    \text{XY chain}
    \quad\longleftrightarrow\quad
    \text{spinless BdG chain}
    \quad\longleftrightarrow\quad
    \text{independent }(k,-k)\text{ two-level systems}.
\end{equation}
The transverse-field Ising chain is the special case $\gamma=1$.  The XX chain is
the special case $\gamma=0$, where the pairing component vanishes and the problem
reduces to a number-conserving free-fermion hopping model.

\begin{remark}[Normalization warning]
Different references place factors of $2$ in different locations in
\cref{eq:bdg-chapter-xy-spin-hamiltonian}.  The invariant statement is the shape
of the spectrum and the gap-closing conditions.  With the above convention, the
critical lines are obtained from
\begin{equation}
    h-\cos k=0,
    \qquad
    \gamma\sin k=0.
\end{equation}
For $\gamma\neq0$, this gives $h=1$ at $k=0$ and $h=-1$ at $k=\pi$.
\end{remark}

\section{Nambu space versus sublattice space}

A two-component Hamiltonian of the form
\begin{equation}
    \mathcal{H}(k)=d_x(k)\tau_x+d_y(k)\tau_y+d_z(k)\tau_z
\end{equation}
may describe very different physics depending on what the two components mean.
For the Kitaev and XY chains, the two components are Nambu components,
\begin{equation}
    \Psi_k=(c_k,c_{-k}^\dagger)^T.
\end{equation}
For the SSH chain, the two components are sublattice components,
\begin{equation}
    \psi_k=(a_k,b_k)^T,
\end{equation}
where $a_k$ and $b_k$ annihilate fermions on two distinct sites within the unit
cell.  The formal $2\times2$ matrix notation is similar, but the physical meaning
is different.

\begin{center}
\begin{tabularx}{0.96\linewidth}{@{}L{0.24\linewidth}YY@{}}
\toprule
 & BdG/Nambu two-level system & SSH/sublattice two-band system \\
\midrule
Two-component spinor
& $(c_k,c_{-k}^\dagger)^T$
& $(a_k,b_k)^T$ \\
Physical components
& particle and hole of related momenta
& two independent orbitals/sublattices in one unit cell \\
Reason for two bands
& artificial Nambu doubling
& genuine Bloch-band structure \\
Built-in symmetry
& particle--hole redundancy
$\tau_x\mathcal{H}(-k)^*\tau_x=-\mathcal{H}(k)$
& no BdG redundancy; SSH has chiral symmetry only for special hopping structure \\
Prefactor in Hamiltonian
& $\frac12\sum_k\Psi_k^\dagger\mathcal{H}_{\rm BdG}(k)\Psi_k$
& $\sum_k\psi_k^\dagger h_{\rm Bloch}(k)\psi_k$ \\
Meaning of negative energy
& hole partner of positive-energy quasiparticle
& physical valence band \\
Ground state
& Bogoliubov/BCS quasiparticle vacuum
& filled-band Slater determinant at fixed particle number \\
\bottomrule
\end{tabularx}
\end{center}

This distinction is central to the organization of these notes.  The Kitaev,
TFIM, and XY chains are solved through BdG pairing blocks.  The SSH model is a
number-conserving free-fermion Bloch problem.  Both can be represented by maps
from momentum space to a vector $\bm{d}(k)$, and both can carry topological
winding numbers in one dimension, but the Hilbert-space interpretation of the
same-looking Pauli matrices is not the same.

\section{Summary}

The BdG formalism rewrites a quadratic parity-preserving fermion Hamiltonian in
Nambu space.  In real space,
\begin{equation}
    \mathbb{H}_{\rm BdG}
    =
    \begin{pmatrix}
        h&\Delta\\
        -\Delta^*&-h^T
    \end{pmatrix},
    \qquad
    h^\dagger=h,
    \qquad
    \Delta^T=-\Delta.
\end{equation}
In a translation-invariant spinless chain, the basic momentum block is
\begin{equation}
    \mathcal{H}_{\rm BdG}(k)
    =
    \begin{pmatrix}
        \xi(k)&\Delta_\phi(k)\\
        \Delta_\phi(k)^*&-\xi(-k)
    \end{pmatrix}.
\end{equation}
For $\xi(-k)=\xi(k)$ this becomes a two-level Hamiltonian
\begin{equation}
    \mathcal{H}_{\rm BdG}(k)=\bm{d}(k)\cdot\bm{\tau},
    \qquad
    \bm{d}(k)=
    \bigl(\operatorname{Re}\Delta_\phi(k),
    -\operatorname{Im}\Delta_\phi(k),
    \xi(k)\bigr),
\end{equation}
with eigenvalues
\begin{equation}
    E_\pm(k)=\pm\sqrt{\xi(k)^2+\abs{\Delta_\phi(k)}^2}.
\end{equation}
The particle--hole relation
\begin{equation}
    \tau_x\mathcal{H}_{\rm BdG}(-k)^*\tau_x=-\mathcal{H}_{\rm BdG}(k)
\end{equation}
is an intrinsic Nambu redundancy.  The next chapter will use this structure to
construct the explicit Bogoliubov transformation, the quasiparticle operators,
and the BCS product ground state.
